\subsection{Line Integrals}

\content{

}{% presentation
\vspace{4in}
}{% nopresentation
}{% comments
Introduce the concept of a line integral. Derive the Reimann Sum 
$$ \sum_{i=1}^n f(x_i^*,y_i^*)\Delta s_i $$
}

\content{
\begin{theorem} If $f$ is continuous, and $C$ is a smooth curve, then 
$$ \int_C f(x,y)\, ds = \int_a^b f(x(t),y(t))\sqrt{x'(t)^2+y'(t)^2}\,dt. $$
\end{theorem}
}{% presentation
}{% nopresentation
}{% comments
}

\content{
\begin{example} Evaluate the integral $\int_C (2+x^2y)\,ds$ where $C$ is the
upper half of the unit circle. 
\end{example}
}{% presentation
\vspace{3in}
}{% nopresentation
}{% comments
Answer: $2\pi+\frac{2}{3}$. 
}

\content{
\begin{example} Evaluate the integral $\int_C 2x\,ds$ where $C$ consists of the
arc $C_1$ of the parabola $y=x^2$ from $(0,0)$ to $(1,1)$ followed by the
vertical line segment $C_2$ from $(1,1)$ to $(1,2)$.
\end{example}
}{% presentation
\vspace{3in}
}{% nopresentation
}{% comments
}

\subsubsection{Integrals with Respect to $x$ or $y$}

\content{
\begin{theorem} The Line integrals with respect to $x$ and $y$ are given by 
$$ \int_C f(x,y)dx = \int_a^b f(x(t),y(t))x'(t)\,dt, $$
and 
$$ \int_C f(x,y)dy = \int_a^b f(x(t),y(t))y'(t)\,dt. $$
\end{theorem}
}{% presentation
\vspace{2in}
}{% nopresentation
}{% comments
We will see that these often occur together, in which case we usually use the
following abbreviation: 
\begin{eqnarray*}\int_C P(x,y)dx+\int_C Q(x,y)dy = \\
\int_C P(x,y)dx+Q(x,y)dy. \end{eqnarray*}
}

\content{
\begin{example} Evaluate the integral $\int_Cy^2dx+xdy$ where $C$ is the line
segment from $(-5,-3)$ to $(0,2)$.
\end{example}
}{% presentation
\vspace{4in}
}{% nopresentation
}{% comments
}

\content{
We have a similar formula for line integrals in space: 
$$ \int_C f(x,y,z)\, ds = \int_a^b f(x(t),y(t),z(t))\sqrt{x'(t)^2+
y'(t)^2+z'(t)^2}\,dt. $$
This (and the formula for line integrals in the plane) can be written in the
more compact form 
$$ \int_C f\, ds = \int_C f(\myvec{r}(t))|\myvec{r}'(t)|\, dt. $$
}{% presentation
}{% nopresentation
}{% comments
}

\content{
\begin{example} Evaluate the line integral $\int_C ydx+zdy+xdz$ where $C$ is the
line segment from $(2,0,0)$ to $(3,4,5)$.
\end{example}
}{% presentation
\vspace{4in}
}{% nopresentation
}{% comments
}

\subsubsection{Line Integrals of Vector Fields and Work}
\content{

}{% presentation
\vspace{3in}
}{% nopresentation
}{% comments
Derive the Riemann Sum for the work done by the force field $\myvec{F}$ in
moving a particle along the path $C$: 
$$\sum_{i=1}^n [\myvec{F}(x_i^*,y_i^*,z_i^*)\cdot
\myvec{T}(x_i^*,y_i^*,z_i^*)]\Delta s_i. $$
Do this by finding the component of $\myvec{F}(x_i^*,y_i^*,z_i^*)$ in the
direction of $\myvec{r}'(t)$. 
}

\content{
\begin{definition} We define the work done by $\myvec{F}$ in moving a particle
along the curve $C$ by 
$$ W = \int_C \myvec{F}(x,y,z)\cdot \myvec{T}(x,y,z)\,ds=\int_C \myvec{F}\cdot
\myvec{T}\,ds.  $$
\end{definition}
}{% presentation
\vspace{2in}
}{% nopresentation
}{% comments
If the curve is parametrized by $\vec{r}(t)$, derive the formula 
$$ W = \int_a^b f(\myvec{r}(t))\cdot \myvec{r}'(t)\,dt $$
which is denoted 
$$\int_C f(\myvec{r}(t))\cdot d\myvec{r}. $$
}

\content{
\begin{example} Find the work done by the force field
$\myvec{F}(x,y) = \la x^2, -xy\ra $ in moving a particle along the
quarter-circle $\myvec{r}(t) = \la \cos t,\sin t\ra$ $0\leq t\leq \pi/2$.
\end{example}
}{% presentation
\vspace{4in}
}{% nopresentation
}{% comments
}
